% !PRINTED 327
\weberpage{327}
%<<CONT
according to the algorithm of the greatest common divisor. For from the congruences
\[
\begin{aligned}
D &\equiv D' \pmod{4n}, \\
n &\equiv n' \pmod{D'}, \\
D' &\equiv D'' \pmod{4n'}, \\
n' &\equiv n'' \pmod{D''}, \\
&\cdots\cdots\cdots\cdots
\end{aligned}
\]
one can determine the sequence of numbers
\[
D', \quad 2n', \quad D'', \quad 2n'', \quad \dots
\]
in such a way that each following one is smaller in absolute value than the preceding one, so long as none of the numbers $n'$, $n''$, $n'''$, $\dots$ has become equal to $\pm 1$.

If $D$ is not a square number, then a number $\beta$ can always be determined so that
\weq{21}{(D, \beta) = -1.}

For if $\varDelta$ is the fundamental discriminant associated with $D$ and $\beta$ is relatively prime to $D$, then (by 9)
\[
(D, \beta) = (\varDelta, \beta).
\]

If $\varDelta = \delta \varDelta'$ and $\delta$ is a prime discriminant, and therefore relatively prime to $\varDelta'$, then $\beta_0$ can be assumed in such a way that
\[
(\delta, \beta_0) = -1
\]
holds [by (10) and (11)], and $\beta$ can be determined from the two congruences
\[
\begin{aligned}
\beta &\equiv \beta_0 \pmod{\delta}, \\
&\equiv 1 \pmod{\varDelta'}
\end{aligned}
\]
Then one obtains, by (12),
\[
(D, \beta) = (\delta, \beta_0) = -1.
\]

If one lets $s$ run through a complete system of residues with respect to the modulus $D$ and sets
\[
S = \overset{s}{\Sigma}(D, s), % Source scan confirms s as the upper index on Sigma.
\]
then multiplication by $(D, \beta)$ gives
\[
-S = \overset{s}{\Sigma}(D, \beta s),
\]
and since $\beta s$ runs through a system of residues at the same time as $s$, it follows that
\weq{22}{S = 0.}

From this it follows:
%CONT>>
